% !TeX root = main.tex
\IfFileExists{../apunte_mecanica_clasica.tex}{%
  \documentclass[../apunte_mecanica_clasica.tex]{subfiles}%
}{%
  \documentclass[apunte_mecanica_clasica.tex]{subfiles}%
}

\begin{document}

\chapter{Introducción}

Todos percibimos un Universo de objetos que se mueven alrededor nuestro, y en consecuencia tratamos de entenderlo para luego posteriormente reproducirlos en un laboratorio o en el papel.

\section{Tres conceptos fundamentales}

Con el fin de cumplir con este objetivo se deben introducir \textbf{tres definiciones básicas}, que son fundamentales de la física:

\begin{itemize}
    \item \textbf{Espacio:} El espacio físico es el lugar que contiene a los objetos, es el escenario en los que ocurren los eventos (sucesos), y además \textit{posee tres direcciones independientes}.
    
    \item \textbf{Tiempo:}  El tiempo físico es una\footnote{Una \textbf{magnitud física} es toda propiedad de un objeto o fenómeno que puede ser medida de manera objetiva y expresada numéricamente, acompañada de una unidad. Ejemplos: masa, longitud, tiempo, temperatura, carga eléctrica.} \textbf{magnitud física} que nos permite medir la duración o separación de los acontecimientos (sucesos), sujetos a cambios de los sistemas bajo observación. Esta medida permite \textit{ordenar los sucesos en secuencias}, definiendo un antes y un después de un evento rigiéndose por el\footnote{El \textbf{principio de causalidad} establece que todo suceso (efecto) tiene una causa que lo antecede. Este principio permite ordenar los eventos en el tiempo y es esencial en la formulación de teorías físicas: ningún efecto puede preceder a su causa.} \textbf{principio de causalidad}.
    
    \item \textbf{Materia:} La materia es la \textit{sustancia de la que los objetos están construidos} y ocupa una región del espacio y tiene permanencia en el tiempo.
\end{itemize}

Estos tres conceptos son definidos de forma independientes. Sin embargo, cuando la materia se mueve, estos tres principios se unifican y dan origen a un Universo\footnote{La \textbf{cinemática} es la rama de la mecánica que estudia el movimiento de los cuerpos sin considerar las causas que lo producen. Se enfoca en describir trayectorias, velocidades y aceleraciones.} \textbf{cinemático} y\footnote{La \textbf{dinámica} es la parte de la mecánica que estudia las causas del movimiento, es decir, las fuerzas que actúan sobre los cuerpos y cómo estas determinan su comportamiento.} \textbf{dinámico}, o sea un Universo\footnote{La \textbf{mecánica} es la rama de la física que estudia el movimiento de los cuerpos y las fuerzas que lo producen. Incluye tanto la cinemática como la dinámica.} \textbf{mecánico}.

\section{La materia en la mecánica clásica}

La definición de materia es una definición muy general, puesto que no especifica de qué están construidos los objetos.

Sin embargo, \textit{en la mecánica clásica solo importa la masa de la materia y no de qué está construida la materia} debido a la universalidad del movimiento (no importa la sustancia de que están hechos los cuerpos), por lo que solo consideraremos la masa del cuerpo:

\begin{apuntebloque}{El rol de la materia en la mecánica}
\textit{``La masa es una medida de la cantidad de materia que posee un objeto.''}
\end{apuntebloque}

\section{La noción del tiempo}

La definición del tiempo que usaremos aquí corresponde a la visión de Newtoniana:
\begin{itemize}
    \item El tiempo es absoluto.
    \item El tiempo en el marco Newtoniana es independiente del espacio y la materia.
\end{itemize} 

\subsection*{Visión aristotélica}

No obstante, casi 1500 años antes de Newton, Aristóteles creía que el espacio, el tiempo y la materia no eran \textit{entes independientes} y estableció:

\begin{apuntebloque}{Aristóteles sobre el tiempo}
\textit{``El tiempo es el número del movimiento, número según lo anterior y lo posterior entre un antes y un después. Ahora bien, es imposible que se generen o destruyan el movimiento, ya que no podrían existir el antes y el después si no hubiera tiempo, y ciertamente, el movimiento es continuo como el tiempo, pues este o es lo mismo o es una afección del movimiento.''}
\end{apuntebloque}

Para Aristóteles sólo existe el movimiento, ni el espacio ni el tiempo y ni la materia existen por sí solas. 

\begin{apuntebloque}{Siglos antes de Aristóteles Heráclito de Éfeso afirmó}
\textit{``El movimiento es la ley de la naturaleza. Todo fluye.''}
\end{apuntebloque}

Esta idea de unificación fue recogida por Einstein en su Teoría General de la Relatividad, que no consideraremos aquí. 

\newpage

\section{Tipos de movimiento}

En este curso consideraremos dos tipos de movimientos:
\begin{itemize}
    \item \textbf{Movimiento Físico:} Es el movimiento que se realiza en la Sustancia Real, en la realidad.
    \item \textbf{Movimiento Geométrico:} Es un movimiento abstracto, que no se realiza en la realidad.
\end{itemize}

\section{Ejemplos de movimiento geométrico}

Un ejemplo de movimiento geométrico es de la geometría básica donde una recta se define como \textbf{El lugar geométrico de un punto en movimiento}, definiendo una dirección: El punto geométrico (sin volumen) se puede mover en una dirección, definida por la misma recta, pero en dos sentido: Sentido $S_1$ y sentido $S_2$, tal que $S_1=-S_2$.

\begin{figure}[ht]
    \centering
    \includegraphics[width=0.3\linewidth]{figures/00_introduccion/mov_geo1.png}
    \caption{El lugar geométrico de un punto en movimiento.}
\end{figure}

Otro ejemplo de un movimiento geométrico es cuando graficamos una función, pues movemos la variable independiente desde un punto inicial a un punto final. Por ejemplo $f(x) = e^x$ con $x \in \mathbb{R}$, el punto $x$ se mueve desde $-\infty$ a $+\infty$, generando un movimiento en el eje $y$. Por lo tanto, hay un movimiento geométrico en $x$ (independiente) e $y$ (dependiente).

\begin{figure}[h]
    \centering
    \includegraphics[width=0.3\linewidth]{figures/00_introduccion/mov_geo2.png}
    \caption{Gráfica de la función $f(x)=e^x$.}
\end{figure}

\section{Lugar, posición y movimiento}

En el estudio del movimiento físico, para definir el movimiento primero debemos definir lo que es el lugar de un cuerpo:

\begin{apuntebloque}{Definición de lugar de un cuerpo}
Un cuerpo del Universo se dice que ocupa un lugar en el espacio, dentro del Universo, con respecto al resto de la materia del Universo.
\end{apuntebloque}

Esta definición nos dice que una porción del Universo mantendrá ciertas distancias, con respecto al resto de la materia del Universo.

Esta definición establece que un cuerpo posee una \textbf{posición} en el espacio, que Aristóteles denominó \textbf{lugar}, relativo al resto de la materia del Universo. Esto significa que la posición de una porción del Universo es conocida cuando se conocen las distancias del cuerpo al resto de la materia del Universo.

A partir de la definición de posición o lugar, podemos definir lo que es movimiento:

\begin{apuntebloque}{Definición de movimiento}
El movimiento es el cambio de lugar de una porción del Universo con respecto al resto de la materia del Universo.
\end{apuntebloque}

Esta definición nos dice que una porción del Universo se mueve si sus distancias relativas con respecto al resto del Universo cambian.

\section{La necesidad de múltiples cuerpos}

Se debe notar que la definición del movimiento menciona a \textbf{toda la materia del Universo}, entonces aparece de forma natural la pregunta \textit{¿Por qué se necesita toda la materia del Universo para definir el movimiento?}
 
La respuesta la podemos entender como sigue: 

\begin{itemize}
    \item Supongamos que sólo existieran dos cuerpos en el Universo, entonces el movimiento circular, como muestra la figura \ref{fig:dos_cuerpos_movimiento_circular}, uno de ellos entorno del otro, no podría ser pensado.
    \begin{figure}[ht]
        \centering
        \includegraphics[width=0.5\linewidth]{figures/00_introduccion/fig1.png}
        \caption{Un cuerpo gira alrededor de otro cuerpo.}
        \label{fig:dos_cuerpos_movimiento_circular}
    \end{figure}

    Los cuerpos siempre tendrán la misma distancia, el radio del círculo, y como no hay un tercer cuerpo no se puede medir un ángulo, por lo tanto, los dos cuerpos no saben que están rotando. \textbf{Por lo tanto, se necesita otro cuerpo}.
\end{itemize}

\clearpage

\begin{itemize}
\item Si ahora existieran sólo tres cuerpos en el Universo, el movimiento mostrado en la figura \ref{fig:tres_cuerpos_movimiento_circular} tampoco podría ser pensado, puesto que las distancias entre los cuerpos $d_{12}$, $d_{13}$ y $d_{23}$ no cambian durante el movimiento. \textbf{Por lo tanto, se necesita otro cuerpo}.

    \begin{figure}[ht]
        \centering
        \includegraphics[width=0.3\linewidth]{figures/00_introduccion/fig2.png}
        \caption{El cuerpo 3 rota en torno al eje de distancia $d_{12}$.}
        \label{fig:tres_cuerpos_movimiento_circular}
    \end{figure}
    Por lo tanto, se necesita un Universo de materia para determinar si un objeto se mueve.
    
\end{itemize}

\textbf{Hacia una física medible:} Ahora que sabemos cómo saber si un cuerpo se mueve, necesitamos establecer algunas reglas básicas para poder realizar mediciones.

%%%%%%%%%%%%%%%%%%%%%%%%%%%%%%%%%%%%%%%%%%%%%%%%%%%%%%%%%%%%%%%%%%%%%%%%%%%%%%%%%%%%%%%%%%%%%%%%%%%%%%%%%%%%%%%%%%%%%%%%%%%%%%%%%%%%%%%%%%%%%%%%%%%%%%%%%%%%%%%%%

\newpage

\end{document}
